\documentclass[../main.tex]{subfiles}
\begin{document}
\chapter{Candidate selection models}
\label{sec:cand_sel}
The series system consists of $m$ components with random lifetimes $\cv_1,\ldots,\cv_m$.
These random variables are \emph{unobservable} and thus, for instance, we cannot know with certainty from a sample of masked system failures which component is the cause of the system failure.
However, we suppose there is some process $A$ that has some knowledge about the components and which generates a set of candidate component causes, as depicted in \cref{fig:true_model}.
\begin{figure}
	\label{fig:true_model}
	\caption{Graphical model of random variables}
	\begin{center}
		% Graphic for TeX using PGF
% Title: /home/spinoza/Documents/repos/papers/series_system_estimation_from_masked_system_failure_time_samples/research/true_graph_model.dia
% Creator: Dia v0.97+git
% CreationDate: Sat Jul  6 06:56:22 2019
% For: spinoza
% \usepackage{tikz}
% The following commands are not supported in PSTricks at present
% We define them conditionally, so when they are implemented,
% this pgf file will use them.
\ifx\du\undefined
  \newlength{\du}
\fi
\setlength{\du}{15\unitlength}
\begin{tikzpicture}[even odd rule]
\pgftransformxscale{1.000000}
\pgftransformyscale{-1.000000}
\definecolor{dialinecolor}{rgb}{0.000000, 0.000000, 0.000000}
\pgfsetstrokecolor{dialinecolor}
\pgfsetstrokeopacity{1.000000}
\definecolor{diafillcolor}{rgb}{1.000000, 1.000000, 1.000000}
\pgfsetfillcolor{diafillcolor}
\pgfsetfillopacity{1.000000}
\pgfsetlinewidth{0.100000\du}
\pgfsetdash{{\pgflinewidth}{0.200000\du}}{0cm}
\pgfsetmiterjoin
\pgfsetbuttcap
\definecolor{dialinecolor}{rgb}{0.498039, 0.498039, 0.498039}
\pgfsetstrokecolor{dialinecolor}
\pgfsetstrokeopacity{1.000000}
\pgfpathmoveto{\pgfpoint{41.000000\du}{10.000000\du}}
\pgfpathcurveto{\pgfpoint{41.000000\du}{10.000000\du}}{\pgfpoint{41.000000\du}{8.000000\du}}{\pgfpoint{44.000000\du}{8.000000\du}}
\pgfpathcurveto{\pgfpoint{47.000000\du}{8.000000\du}}{\pgfpoint{48.000000\du}{8.000000\du}}{\pgfpoint{49.000000\du}{8.000000\du}}
\pgfpathcurveto{\pgfpoint{50.000000\du}{8.000000\du}}{\pgfpoint{51.000000\du}{8.000000\du}}{\pgfpoint{54.000000\du}{8.000000\du}}
\pgfpathcurveto{\pgfpoint{57.000000\du}{8.000000\du}}{\pgfpoint{57.000000\du}{10.000000\du}}{\pgfpoint{57.000000\du}{10.000000\du}}
\pgfpathcurveto{\pgfpoint{57.000000\du}{10.000000\du}}{\pgfpoint{57.000000\du}{12.000000\du}}{\pgfpoint{54.000000\du}{12.000000\du}}
\pgfpathcurveto{\pgfpoint{51.000000\du}{12.000000\du}}{\pgfpoint{50.000000\du}{12.000000\du}}{\pgfpoint{49.000000\du}{12.000000\du}}
\pgfpathcurveto{\pgfpoint{48.000000\du}{12.000000\du}}{\pgfpoint{47.000000\du}{12.000000\du}}{\pgfpoint{44.000000\du}{12.000000\du}}
\pgfpathcurveto{\pgfpoint{41.000000\du}{12.000000\du}}{\pgfpoint{41.000000\du}{10.000000\du}}{\pgfpoint{41.000000\du}{10.000000\du}}
\pgfpathclose
\pgfusepath{stroke}
% setfont left to latex
\definecolor{dialinecolor}{rgb}{0.000000, 0.000000, 0.000000}
\pgfsetstrokecolor{dialinecolor}
\pgfsetstrokeopacity{1.000000}
\definecolor{diafillcolor}{rgb}{0.000000, 0.000000, 0.000000}
\pgfsetfillcolor{diafillcolor}
\pgfsetfillopacity{1.000000}
\node[anchor=base west,inner sep=0pt,outer sep=0pt,color=dialinecolor] at (43.000000\du,10.000000\du){};
% setfont left to latex
\definecolor{dialinecolor}{rgb}{0.000000, 0.000000, 0.000000}
\pgfsetstrokecolor{dialinecolor}
\pgfsetstrokeopacity{1.000000}
\definecolor{diafillcolor}{rgb}{0.000000, 0.000000, 0.000000}
\pgfsetfillcolor{diafillcolor}
\pgfsetfillopacity{1.000000}
\node[anchor=base west,inner sep=0pt,outer sep=0pt,color=dialinecolor] at (43.000000\du,10.000000\du){};
% setfont left to latex
\definecolor{dialinecolor}{rgb}{0.000000, 0.000000, 0.000000}
\pgfsetstrokecolor{dialinecolor}
\pgfsetstrokeopacity{1.000000}
\definecolor{diafillcolor}{rgb}{0.000000, 0.000000, 0.000000}
\pgfsetfillcolor{diafillcolor}
\pgfsetfillopacity{1.000000}
\node[anchor=base west,inner sep=0pt,outer sep=0pt,color=dialinecolor] at (45.013631\du,4.000000\du){};
\pgfsetlinewidth{0.100000\du}
\pgfsetdash{}{0pt}
\pgfsetbuttcap
\pgfsetmiterjoin
\pgfsetlinewidth{0.100000\du}
\pgfsetbuttcap
\pgfsetmiterjoin
\pgfsetdash{}{0pt}
\definecolor{dialinecolor}{rgb}{0.000000, 0.639216, 0.000000}
\pgfsetstrokecolor{dialinecolor}
\pgfsetstrokeopacity{1.000000}
\pgfpathmoveto{\pgfpoint{45.013631\du}{3.000000\du}}
\pgfpathcurveto{\pgfpoint{45.563631\du}{3.000000\du}}{\pgfpoint{46.013631\du}{3.450000\du}}{\pgfpoint{46.013631\du}{4.000000\du}}
\pgfpathcurveto{\pgfpoint{46.013631\du}{4.550000\du}}{\pgfpoint{45.563631\du}{5.000000\du}}{\pgfpoint{45.013631\du}{5.000000\du}}
\pgfpathcurveto{\pgfpoint{44.463631\du}{5.000000\du}}{\pgfpoint{44.013631\du}{4.550000\du}}{\pgfpoint{44.013631\du}{4.000000\du}}
\pgfpathcurveto{\pgfpoint{44.013631\du}{3.450000\du}}{\pgfpoint{44.463631\du}{3.000000\du}}{\pgfpoint{45.013631\du}{3.000000\du}}
\pgfusepath{stroke}
% setfont left to latex
\definecolor{dialinecolor}{rgb}{0.000000, 0.000000, 0.000000}
\pgfsetstrokecolor{dialinecolor}
\pgfsetstrokeopacity{1.000000}
\definecolor{diafillcolor}{rgb}{0.000000, 0.000000, 0.000000}
\pgfsetfillcolor{diafillcolor}
\pgfsetfillopacity{1.000000}
\node[anchor=base west,inner sep=0pt,outer sep=0pt,color=dialinecolor] at (44.739031\du,4.371370\du){S};
% setfont left to latex
\definecolor{dialinecolor}{rgb}{0.000000, 0.000000, 0.000000}
\pgfsetstrokecolor{dialinecolor}
\pgfsetstrokeopacity{1.000000}
\definecolor{diafillcolor}{rgb}{0.000000, 0.000000, 0.000000}
\pgfsetfillcolor{diafillcolor}
\pgfsetfillopacity{1.000000}
\node[anchor=base west,inner sep=0pt,outer sep=0pt,color=dialinecolor] at (49.000000\du,4.000000\du){};
% setfont left to latex
\definecolor{dialinecolor}{rgb}{0.000000, 0.000000, 0.000000}
\pgfsetstrokecolor{dialinecolor}
\pgfsetstrokeopacity{1.000000}
\definecolor{diafillcolor}{rgb}{0.000000, 0.000000, 0.000000}
\pgfsetfillcolor{diafillcolor}
\pgfsetfillopacity{1.000000}
\node[anchor=base west,inner sep=0pt,outer sep=0pt,color=dialinecolor] at (49.000000\du,4.000000\du){};
% setfont left to latex
\definecolor{dialinecolor}{rgb}{0.000000, 0.000000, 0.000000}
\pgfsetstrokecolor{dialinecolor}
\pgfsetstrokeopacity{1.000000}
\definecolor{diafillcolor}{rgb}{0.000000, 0.000000, 0.000000}
\pgfsetfillcolor{diafillcolor}
\pgfsetfillopacity{1.000000}
\node[anchor=base west,inner sep=0pt,outer sep=0pt,color=dialinecolor] at (49.000000\du,4.000000\du){};
\pgfsetlinewidth{0.100000\du}
\pgfsetdash{}{0pt}
\pgfsetbuttcap
\pgfsetmiterjoin
\pgfsetlinewidth{0.100000\du}
\pgfsetbuttcap
\pgfsetmiterjoin
\pgfsetdash{}{0pt}
\definecolor{dialinecolor}{rgb}{0.000000, 0.639216, 0.000000}
\pgfsetstrokecolor{dialinecolor}
\pgfsetstrokeopacity{1.000000}
\pgfpathmoveto{\pgfpoint{49.013631\du}{3.000000\du}}
\pgfpathcurveto{\pgfpoint{49.563631\du}{3.000000\du}}{\pgfpoint{50.013631\du}{3.450000\du}}{\pgfpoint{50.013631\du}{4.000000\du}}
\pgfpathcurveto{\pgfpoint{50.013631\du}{4.550000\du}}{\pgfpoint{49.563631\du}{5.000000\du}}{\pgfpoint{49.013631\du}{5.000000\du}}
\pgfpathcurveto{\pgfpoint{48.463631\du}{5.000000\du}}{\pgfpoint{48.013631\du}{4.550000\du}}{\pgfpoint{48.013631\du}{4.000000\du}}
\pgfpathcurveto{\pgfpoint{48.013631\du}{3.450000\du}}{\pgfpoint{48.463631\du}{3.000000\du}}{\pgfpoint{49.013631\du}{3.000000\du}}
\pgfusepath{stroke}
% setfont left to latex
\definecolor{dialinecolor}{rgb}{0.000000, 0.000000, 0.000000}
\pgfsetstrokecolor{dialinecolor}
\pgfsetstrokeopacity{1.000000}
\definecolor{diafillcolor}{rgb}{0.000000, 0.000000, 0.000000}
\pgfsetfillcolor{diafillcolor}
\pgfsetfillopacity{1.000000}
\node[anchor=base west,inner sep=0pt,outer sep=0pt,color=dialinecolor] at (48.602431\du,4.350340\du){C};
% setfont left to latex
\definecolor{dialinecolor}{rgb}{0.000000, 0.000000, 0.000000}
\pgfsetstrokecolor{dialinecolor}
\pgfsetstrokeopacity{1.000000}
\definecolor{diafillcolor}{rgb}{0.000000, 0.000000, 0.000000}
\pgfsetfillcolor{diafillcolor}
\pgfsetfillopacity{1.000000}
\node[anchor=base west,inner sep=0pt,outer sep=0pt,color=dialinecolor] at (50.000000\du,4.000000\du){};
% setfont left to latex
\definecolor{dialinecolor}{rgb}{0.000000, 0.000000, 0.000000}
\pgfsetstrokecolor{dialinecolor}
\pgfsetstrokeopacity{1.000000}
\definecolor{diafillcolor}{rgb}{0.000000, 0.000000, 0.000000}
\pgfsetfillcolor{diafillcolor}
\pgfsetfillopacity{1.000000}
\node[anchor=base west,inner sep=0pt,outer sep=0pt,color=dialinecolor] at (53.000000\du,4.000000\du){};
% setfont left to latex
\definecolor{dialinecolor}{rgb}{0.000000, 0.000000, 0.000000}
\pgfsetstrokecolor{dialinecolor}
\pgfsetstrokeopacity{1.000000}
\definecolor{diafillcolor}{rgb}{0.000000, 0.000000, 0.000000}
\pgfsetfillcolor{diafillcolor}
\pgfsetfillopacity{1.000000}
\node[anchor=base west,inner sep=0pt,outer sep=0pt,color=dialinecolor] at (54.400000\du,10.400000\du){Tm};
% setfont left to latex
\definecolor{dialinecolor}{rgb}{0.000000, 0.000000, 0.000000}
\pgfsetstrokecolor{dialinecolor}
\pgfsetstrokeopacity{1.000000}
\definecolor{diafillcolor}{rgb}{0.000000, 0.000000, 0.000000}
\pgfsetfillcolor{diafillcolor}
\pgfsetfillopacity{1.000000}
\node[anchor=base west,inner sep=0pt,outer sep=0pt,color=dialinecolor] at (55.000000\du,10.000000\du){};
\pgfsetlinewidth{0.100000\du}
\pgfsetdash{{\pgflinewidth}{0.200000\du}}{0cm}
\pgfsetmiterjoin
\pgfsetbuttcap
{
\definecolor{diafillcolor}{rgb}{0.000000, 0.000000, 0.000000}
\pgfsetfillcolor{diafillcolor}
\pgfsetfillopacity{1.000000}
% was here!!!
\pgfsetarrowsend{latex}
\definecolor{dialinecolor}{rgb}{0.000000, 0.000000, 0.000000}
\pgfsetstrokecolor{dialinecolor}
\pgfsetstrokeopacity{1.000000}
\pgfpathmoveto{\pgfpoint{49.000000\du}{8.000000\du}}
\pgfpathcurveto{\pgfpoint{49.000000\du}{6.000000\du}}{\pgfpoint{45.013631\du}{7.150000\du}}{\pgfpoint{45.013631\du}{5.150000\du}}
\pgfusepath{stroke}
}
\pgfsetlinewidth{0.100000\du}
\pgfsetdash{}{0pt}
\pgfsetbuttcap
\pgfsetmiterjoin
\pgfsetlinewidth{0.100000\du}
\pgfsetbuttcap
\pgfsetmiterjoin
\pgfsetdash{}{0pt}
\definecolor{dialinecolor}{rgb}{0.803922, 0.000000, 0.000000}
\pgfsetstrokecolor{dialinecolor}
\pgfsetstrokeopacity{1.000000}
\pgfpathellipse{\pgfpoint{46.000000\du}{10.000000\du}}{\pgfpoint{1.000000\du}{0\du}}{\pgfpoint{0\du}{1.000000\du}}
\pgfusepath{stroke}
% setfont left to latex
\definecolor{dialinecolor}{rgb}{0.000000, 0.000000, 0.000000}
\pgfsetstrokecolor{dialinecolor}
\pgfsetstrokeopacity{1.000000}
\definecolor{diafillcolor}{rgb}{0.000000, 0.000000, 0.000000}
\pgfsetfillcolor{diafillcolor}
\pgfsetfillopacity{1.000000}
\node[anchor=base west,inner sep=0pt,outer sep=0pt,color=dialinecolor] at (45.400000\du,10.400000\du){T2};
% setfont left to latex
\definecolor{dialinecolor}{rgb}{0.000000, 0.000000, 0.000000}
\pgfsetstrokecolor{dialinecolor}
\pgfsetstrokeopacity{1.000000}
\definecolor{diafillcolor}{rgb}{0.000000, 0.000000, 0.000000}
\pgfsetfillcolor{diafillcolor}
\pgfsetfillopacity{1.000000}
\node[anchor=base west,inner sep=0pt,outer sep=0pt,color=dialinecolor] at (46.000000\du,10.000000\du){};
\pgfsetlinewidth{0.100000\du}
\pgfsetdash{}{0pt}
\pgfsetbuttcap
\pgfsetmiterjoin
\pgfsetlinewidth{0.100000\du}
\pgfsetbuttcap
\pgfsetmiterjoin
\pgfsetdash{}{0pt}
\definecolor{dialinecolor}{rgb}{0.803922, 0.000000, 0.000000}
\pgfsetstrokecolor{dialinecolor}
\pgfsetstrokeopacity{1.000000}
\pgfpathellipse{\pgfpoint{43.000000\du}{10.000000\du}}{\pgfpoint{1.000000\du}{0\du}}{\pgfpoint{0\du}{1.000000\du}}
\pgfusepath{stroke}
% setfont left to latex
\definecolor{dialinecolor}{rgb}{0.000000, 0.000000, 0.000000}
\pgfsetstrokecolor{dialinecolor}
\pgfsetstrokeopacity{1.000000}
\definecolor{diafillcolor}{rgb}{0.000000, 0.000000, 0.000000}
\pgfsetfillcolor{diafillcolor}
\pgfsetfillopacity{1.000000}
\node[anchor=base west,inner sep=0pt,outer sep=0pt,color=dialinecolor] at (42.400000\du,10.400000\du){T1};
\pgfsetlinewidth{0.100000\du}
\pgfsetdash{}{0pt}
\pgfsetbuttcap
\pgfsetmiterjoin
\pgfsetlinewidth{0.100000\du}
\pgfsetbuttcap
\pgfsetmiterjoin
\pgfsetdash{}{0pt}
\definecolor{dialinecolor}{rgb}{0.803922, 0.000000, 0.000000}
\pgfsetstrokecolor{dialinecolor}
\pgfsetstrokeopacity{1.000000}
\pgfpathellipse{\pgfpoint{55.000000\du}{10.000000\du}}{\pgfpoint{1.000000\du}{0\du}}{\pgfpoint{0\du}{1.000000\du}}
\pgfusepath{stroke}
% setfont left to latex
\definecolor{dialinecolor}{rgb}{0.000000, 0.000000, 0.000000}
\pgfsetstrokecolor{dialinecolor}
\pgfsetstrokeopacity{1.000000}
\definecolor{diafillcolor}{rgb}{0.000000, 0.000000, 0.000000}
\pgfsetfillcolor{diafillcolor}
\pgfsetfillopacity{1.000000}
\node[anchor=base west,inner sep=0pt,outer sep=0pt,color=dialinecolor] at (55.000000\du,10.000000\du){};
\pgfsetlinewidth{0.100000\du}
\pgfsetdash{{\pgflinewidth}{0.200000\du}}{0cm}
\pgfsetbuttcap
{
\definecolor{diafillcolor}{rgb}{0.000000, 0.000000, 0.000000}
\pgfsetfillcolor{diafillcolor}
\pgfsetfillopacity{1.000000}
% was here!!!
\definecolor{dialinecolor}{rgb}{0.000000, 0.000000, 0.000000}
\pgfsetstrokecolor{dialinecolor}
\pgfsetstrokeopacity{1.000000}
\draw (47.500000\du,10.000000\du)--(53.451200\du,10.000000\du);
}
% setfont left to latex
\definecolor{dialinecolor}{rgb}{0.000000, 0.000000, 0.000000}
\pgfsetstrokecolor{dialinecolor}
\pgfsetstrokeopacity{1.000000}
\definecolor{diafillcolor}{rgb}{0.000000, 0.000000, 0.000000}
\pgfsetfillcolor{diafillcolor}
\pgfsetfillopacity{1.000000}
\node[anchor=base west,inner sep=0pt,outer sep=0pt,color=dialinecolor] at (53.000000\du,4.000000\du){};
% setfont left to latex
\definecolor{dialinecolor}{rgb}{0.000000, 0.000000, 0.000000}
\pgfsetstrokecolor{dialinecolor}
\pgfsetstrokeopacity{1.000000}
\definecolor{diafillcolor}{rgb}{0.000000, 0.000000, 0.000000}
\pgfsetfillcolor{diafillcolor}
\pgfsetfillopacity{1.000000}
\node[anchor=base west,inner sep=0pt,outer sep=0pt,color=dialinecolor] at (52.842360\du,3.886212\du){};
\pgfsetlinewidth{0.100000\du}
\pgfsetdash{}{0pt}
\pgfsetbuttcap
\pgfsetmiterjoin
\pgfsetlinewidth{0.100000\du}
\pgfsetbuttcap
\pgfsetmiterjoin
\pgfsetdash{}{0pt}
\definecolor{dialinecolor}{rgb}{0.803922, 0.000000, 0.000000}
\pgfsetstrokecolor{dialinecolor}
\pgfsetstrokeopacity{1.000000}
\pgfpathellipse{\pgfpoint{53.007601\du}{4.000000\du}}{\pgfpoint{1.000000\du}{0\du}}{\pgfpoint{0\du}{1.000000\du}}
\pgfusepath{stroke}
% setfont left to latex
\definecolor{dialinecolor}{rgb}{0.000000, 0.000000, 0.000000}
\pgfsetstrokecolor{dialinecolor}
\pgfsetstrokeopacity{1.000000}
\definecolor{diafillcolor}{rgb}{0.000000, 0.000000, 0.000000}
\pgfsetfillcolor{diafillcolor}
\pgfsetfillopacity{1.000000}
\node[anchor=base west,inner sep=0pt,outer sep=0pt,color=dialinecolor] at (52.596401\du,4.350340\du){K};
% setfont left to latex
\definecolor{dialinecolor}{rgb}{0.000000, 0.000000, 0.000000}
\pgfsetstrokecolor{dialinecolor}
\pgfsetstrokeopacity{1.000000}
\definecolor{diafillcolor}{rgb}{0.000000, 0.000000, 0.000000}
\pgfsetfillcolor{diafillcolor}
\pgfsetfillopacity{1.000000}
\node[anchor=base west,inner sep=0pt,outer sep=0pt,color=dialinecolor] at (52.787834\du,4.036156\du){};
% setfont left to latex
\definecolor{dialinecolor}{rgb}{0.000000, 0.000000, 0.000000}
\pgfsetstrokecolor{dialinecolor}
\pgfsetstrokeopacity{1.000000}
\definecolor{diafillcolor}{rgb}{0.000000, 0.000000, 0.000000}
\pgfsetfillcolor{diafillcolor}
\pgfsetfillopacity{1.000000}
\node[anchor=base west,inner sep=0pt,outer sep=0pt,color=dialinecolor] at (52.787834\du,3.899843\du){};
\pgfsetlinewidth{0.100000\du}
\pgfsetdash{{\pgflinewidth}{0.200000\du}}{0cm}
\pgfsetmiterjoin
\pgfsetbuttcap
{
\definecolor{diafillcolor}{rgb}{0.000000, 0.000000, 0.000000}
\pgfsetfillcolor{diafillcolor}
\pgfsetfillopacity{1.000000}
% was here!!!
\pgfsetarrowsend{latex}
\definecolor{dialinecolor}{rgb}{0.000000, 0.000000, 0.000000}
\pgfsetstrokecolor{dialinecolor}
\pgfsetstrokeopacity{1.000000}
\pgfpathmoveto{\pgfpoint{49.000000\du}{8.000000\du}}
\pgfpathcurveto{\pgfpoint{49.000000\du}{6.000000\du}}{\pgfpoint{49.000000\du}{7.150000\du}}{\pgfpoint{49.000000\du}{5.150000\du}}
\pgfusepath{stroke}
}
\pgfsetlinewidth{0.100000\du}
\pgfsetdash{{\pgflinewidth}{0.200000\du}}{0cm}
\pgfsetmiterjoin
\pgfsetbuttcap
{
\definecolor{diafillcolor}{rgb}{0.000000, 0.000000, 0.000000}
\pgfsetfillcolor{diafillcolor}
\pgfsetfillopacity{1.000000}
% was here!!!
\pgfsetarrowsend{latex}
\definecolor{dialinecolor}{rgb}{0.000000, 0.000000, 0.000000}
\pgfsetstrokecolor{dialinecolor}
\pgfsetstrokeopacity{1.000000}
\pgfpathmoveto{\pgfpoint{49.000000\du}{8.000000\du}}
\pgfpathcurveto{\pgfpoint{49.000000\du}{6.000000\du}}{\pgfpoint{53.007601\du}{7.150000\du}}{\pgfpoint{53.007601\du}{5.150000\du}}
\pgfusepath{stroke}
}
% setfont left to latex
\definecolor{dialinecolor}{rgb}{0.000000, 0.000000, 0.000000}
\pgfsetstrokecolor{dialinecolor}
\pgfsetstrokeopacity{1.000000}
\definecolor{diafillcolor}{rgb}{0.000000, 0.000000, 0.000000}
\pgfsetfillcolor{diafillcolor}
\pgfsetfillopacity{1.000000}
\node[anchor=base west,inner sep=0pt,outer sep=0pt,color=dialinecolor] at (47.049059\du,5.576493\du){A};
% setfont left to latex
\definecolor{dialinecolor}{rgb}{0.000000, 0.000000, 0.000000}
\pgfsetstrokecolor{dialinecolor}
\pgfsetstrokeopacity{1.000000}
\definecolor{diafillcolor}{rgb}{0.000000, 0.000000, 0.000000}
\pgfsetfillcolor{diafillcolor}
\pgfsetfillopacity{1.000000}
\node[anchor=base west,inner sep=0pt,outer sep=0pt,color=dialinecolor] at (50.375095\du,5.440180\du){};
\pgfsetlinewidth{0.050000\du}
\pgfsetdash{{\pgflinewidth}{0.200000\du}}{0cm}
\pgfsetmiterjoin
\pgfsetbuttcap
{
\definecolor{diafillcolor}{rgb}{0.313726, 0.313726, 0.313726}
\pgfsetfillcolor{diafillcolor}
\pgfsetfillopacity{1.000000}
% was here!!!
\pgfsetarrowsend{stealth}
\definecolor{dialinecolor}{rgb}{0.313726, 0.313726, 0.313726}
\pgfsetstrokecolor{dialinecolor}
\pgfsetstrokeopacity{1.000000}
\pgfpathmoveto{\pgfpoint{47.663599\du}{5.395892\du}}
\pgfpathcurveto{\pgfpoint{48.658683\du}{5.395892\du}}{\pgfpoint{47.895331\du}{6.200138\du}}{\pgfpoint{48.985834\du}{6.241032\du}}
\pgfusepath{stroke}
}
\end{tikzpicture}

	\end{center}
\end{figure}


We consider two ways failed component candidate sets may be constructed.
Suppose candidate sets are of cardinality $w$.
The first way candidate sets are generated is \emph{optimal} since we simply draw from the component lifetimes and include the $w$ components with the smallest lifetimes in the set. The probability that a component in a set $\Set{A} \in \PS{\Set{K}}$ is the cause of a system failure follows from the definition of $\RV{K}$, i.e., $\Prob{\RV{K} \in \Set{A}}$.

The second way is perhaps more realistic in that it is \emph{certain} that one of the components in the candidate set is in a failed state (and is thus the cause of the system failure), e.g., a diagnostician was able to isolate the problem to some subset of the system.
However, the $w-1$ non-failed componets in the candidate set are \emph{randomly} selected from the remaining $m-1$ components, which is probably overly pessimistic.
These two approaches generate a range, from useful but pessimistic to unreasonably \emph{optimal}. For now, we provide notation for the \emph{random component sets} and treat the two approaches separately in \cref{sec:opt,sec:pess}.


The \emph{cardinality} of a component set could theoretically convey 
information about the system. As a simplification, we make the following 
assumption.
\begin{assumption}
	\label{asm:indep_W_S}
	The cardinality of the random set $\rvcand$, denoted by the discrete random variable
	\begin{equation}
	\RV{W} = \Card{\rvcand}\,,
	\end{equation}
	is independent of the random system lifetime $\sv$ and the random component 
	cause $\rvcand$ and therefore does not carry information about the true 
	parameter index $\tparam\sysparam$.
\end{assumption}

The indicator function and the $k$-subset functions are needed for subsequent material and are given by the following definitions.
\begin{definition}[Indicator function]
	Given a subset $\Set{A}$ of a set $\Set{B}$, the indicator function $\SetIndicator{\Set{A}} \colon \Set{B} \mapsto \{0,1\}$
	is equal to $1$ if $x \in \Set{A}$ and $0$ otherwise.
\end{definition}

\begin{definition}
	The subsets of $\Set{A}$ that are of cardinality $k$ is defined as
	\begin{equation}
	\left[\Set{A}\right]^k \coloneqq \SetBuilder{\Set{X} \in \Set{A}}{\Card{\Set{A}} = k}\,,
	\end{equation}
	which has a cardinality
	\begin{equation}
	\binom{\Card{\Set{A}}}{k}\,.
	\end{equation}
\end{definition}

We are primarily interested in the distribution of candidate sets on the condition that they have a specified cardinality.
\begin{definition}
	The random variable $\rvcand$ conditioned on $\RV{W} = w$ has a sample space given by
	\begin{equation}
	\candsw{w} = \left[\PS{\RV{K}}\right]^w\,.
	\end{equation}
\end{definition}

According to \cref{def:random_candidate_set}, the random candidate set $\rvcand$ is a subset of components that are in some way considered to be a plausible cause of system failure.
By \cref{asm:indep_W_S}, the distribution of $\RV{W}$ is not dependent on the true parameter index $\tparam\sysparam$.
\begin{definition}
\label{def:pmf_w}
The probability that $\RV{W} = \Card{\rvcand} = w$ is given by the probability mass function $\pmf_{\RV{W}}\left(w\right)$.
\end{definition}
Any inferences about the parametric model will be \emph{independent} of $\RV{W}$ except in the sense that conditioned on larger candidate sets generates estimators with greater variance.
If a sample of masked system failure times has variable sized candidate sets, we are free to use any parameter-free estimator of the probability mass function of $\RV{W}$, e.g., the empirical distribution as described in \cref{appendix:generalized_general_joint_likelihood_cand}.

By \cref{asm:indep_W_S}, the random lifetime $\sv$ and the cardinality of the random candidate set $\RV{W} = \Card{\rvcand}$ are independent, thus by the axioms of probability the joint density of $\rvcand$, $\sv$, and $\RV{W}$ is given by
\begin{equation}
\label{eq:general_joint_w_cand_likelihood}
\pdf_{\rvcand, \sv, \RV{W}}(\cand, t, w \given \tparam\sysparam) = \pmf_{\rvcand \given \sv, \RV{W}}(\cand \given t, w, \tparam\sysparam) \pmf_{\RV{W}}(w) \spdf(t \given \tparam\sysparam)\,.
\end{equation}
If it is given that $\RV{W} = w$, by the axioms of probability
\begin{equation}
\label{eq:general_joint_cond_w_cand_likelihood}
\jointpdf(\cand, t \given w, \tparam\sysparam) =
\pmf_{\rvcand \given \sv, \RV{W}}(\cand \given t, w, \tparam\sysparam) \spdf(t \given \tparam\sysparam)\,.
\end{equation}


\section{Failed $\alpha$-masked candidate sets}
We consider failed candidate sets given by the following definition.
\begin{definition}
The $\alpha$-mask is a failed candidate set in which the probability the failed component is included in the set is $\alpha \cdot 100 \%$.
\end{definition}

The $\alpha$-mask has a conditional probability distribution given by the following theorem.
\begin{theorem}
The $\alpha$-mask $\rvcand_\alpha$ has a conditional probability distribution given by
\begin{equation}
\pdf_{\rvcand_\alpha \given \RV{K}, \RV{W}}(\cand \given k, w, \alpha) = \frac{\alpha}{\binom{m-1}{w-1}} \SetIndicator{\cand}(k) + \frac{1-\alpha}{\binom{m-1}{w}} \SetIndicator{\SetComplement[\cand]}(k)\,.
\end{equation}
\end{theorem}
\begin{proof}
	Proof here.
	
	We uniformly choose among sets satisfying the criteria.
\end{proof}


The joint distribution is given by
\begin{equation}
	\pdf_{\rvcand,\RV{K},\sv \given \RV{W}}(\cand,k,t \given w, \alpha) = \pdf_{\rvcand \given \RV{K}, \RV{W}}(\cand \given k, w, \alpha) \pdf_{\RV{K},\sv}(k,t \given \tparam\sysparam) \,.
\end{equation}
However, given only a masked failure set, $\RV{K}$ is \emph{unobservable}, and so we \emph{marginalize} over $\RV{K}$.
\begin{theorem}
\begin{equation}
\pdf_{\rvcand,\sv \given \RV{W}}(\cand,t \given w, \tparam\sysparam) = \sum_{k=1}^{m} \pdf_{\rvcand,\RV{K},\sv \given \RV{W}}(\cand,k,t \given w, \alpha)\,.
\end{equation}
\end{theorem}
%which results in
%\begin{equation}
%\pdf_{\rvcand,\sv \given \RV{W}}(\cand,t \given w, \tparam\sysparam) = \srv_{\sv}(t\given \tparam\sysparam) \sum_{k=1}^{m}
%	\haz_k(t \given \tparam\sysparam)
%	\left[
%		\frac{\alpha}{\binom{m-1}{w-1}} \SetIndicator{\cand}(k) + \frac{1-\alpha}{\binom{m-1}{w}} \SetIndicator{\SetComplement[\cand]}(k)	
%	\right] \SetIndicator{\candsw{w}}(\cand)\,.
%\end{equation}

\section{\emph{Optimal} failed candidate sets}
\label{sec:opt}
Optimal candidate sets are generated by drawing a candidate set $\cand$ of cardinality $w$ with a probability proportional to $\Prob{\RV{K} \in \cand}$.
However, the total probability over $\candsw{w}$ does not sum to unity since component sets are not disjoint events with respect to the sample space of $\RV{K}$.
We consider another random element (random set) in which the sample space is defined by the outcomes in $\candsw{w}$.



The joint probability density of $\rvcand$ and $\sv$ conditioned on 
$\RV{W} = w$ is given by the following theorem.
\begin{theorem}
\label{thm:general_joint_likelihood_cand}
The joint probability density of $\rvcand$ and $\sv$ conditioned on $\RV{W} = 
w$ is given by
\begin{equation}
\label{eq:general_joint_likelihood_cand}
    \jointpdf(\cand, t \given w, \tparam\sysparam) = \eta
        \srv_{\sv}\left(t \given \tparam\sysparam\right)
        \sum_{k \in \cand} \haz_k\left(t \given \tparam\sysparam\right)
        \SetIndicator{\candsw{w}}(\cand)\,.
\end{equation}
where $\eta$ is the normalizing constant that makes the joint probability density function integrate\footnote{Integration refers to both summations and integrations.} to unity over the sample space $\candsw{w} \times \left(0,\infty\right)$.
\end{theorem}
\begin{proof}
Given a candidate set $\cand \in \candsw{w}$, the event $E$ we wish to compute the likelihood of is given by
\begin{equation}
    E = \bigcup_{j \in \cand} \left ( \RV{K} = j \cap \sv = t \right )\,.
\end{equation}
We would like to compute the likelihood $L(E)$. The likelihood that $\RV{K} = j$ and $\sv = t$ is given by the probability density function
\begin{equation*}
    \pdf_{\RV{K}, \sv}\left(k, t\given \tparam\sysparam\right) = \haz_k\left(t \given \tparam{\sysparam}_k\right) \srv_{\sv}\left(t \given \tparam\sysparam\right)\,.
    \tag{\ref{eq:k_s_haz} revisited}
\end{equation*}
Performing this substitution yields
\begin{equation}
    L(E) = \bigcup_{j \in \cand} \haz_j\left(t \given \tparam{\sysparam}_j\right) \srv_{\sv}\left(t \given \tparam\sysparam\right)\,.
\end{equation}
Since the events $\RV{K} = j, \sv = t$ for $j \in \cand$ are mutually exclusive, the likelihood is given by
\begin{equation}
    L(E) = \sum_{j \in \cand} \haz_j\left(t \given \tparam{\sysparam}_j\right) \srv_{\sv}\left(t \given \tparam\sysparam\right)\,.
\end{equation}
Since outcomes in the sample space $\candsw{w}$ are not atomic with respect to 
the sample space of $\RV{K} \in \Set{K}$, we multiply by a normalizing constant 
$\eta$ to make it a proper joint probability density function.
\end{proof}

The normalizing constant $\eta$ indicated in \cref{eq:general_joint_likelihood_cand} is given by the following theorem.
\begin{theorem}
The normalizing constant $\eta$ that makes the joint probability density $\jointpdf(\cand, t \given w, \tparam\sysparam)$ integrate to unity over the sample space $\candsw{w} \times \left(0,\infty\right)$ is given by
\begin{equation}
    \eta = \frac{1}{\binom{m-1}{w-1}}\,.
\end{equation}
\end{theorem}
\begin{proof}
Integrating over the sample space $\candsw{w} \times \left(0,\infty\right)$ must be equal to unity. That is,
\begin{equation}
    \eta \sum_{\cand \in \candsw{w}} \int_{0}^{\infty} \jointpdf(\cand, t \given \tparam\sysparam) \dif t = 1\,.
\end{equation}
Substituting the joint probability density with the definition given by \cref{eq:general_joint_likelihood_cand} results in
\begin{equation}
    \eta \sum_{\cand \in \candsw{w}} \int_{0}^{\infty} \left[ \sum_{j \in \cand} \pdf_{\RV{K}, \sv}(j, t \given \tparam\sysparam) \right] \dif t = 1\,.
\end{equation}
We can switch the summation over $\candsw{w}$ and the integration over $t$ without changing the equation, resulting in
\begin{equation}
    \eta \sum_{\cand \in \candsw{w}} \sum_{j \in \cand} \left[ \int_{0}^{\infty} \pdf_{\RV{K}, \sv}(j, t \given \tparam\sysparam) \mathrm{d}t \right] = 1\,.
\end{equation}
By the axioms of probability, the inner integral is the marginal probability mass function of $\RV{K}$. Performing this substitution results in
\begin{equation}
    \eta \sum_{\cand \in \candsw{w}} \sum_{j \in \cand} \pmf_{\RV{K}}(j \given \tparam\sysparam) = 1\,.
\end{equation}
When constructing a candidate set of cardinality $w$, suppose we restrict our attention to only those sets in which component $r$ is present. Then, we have $w-1$ additional components to choose from of the $m-1$ remaining. There are $\binom{m-1}{w-1}$ such choices possible, thus component $r$ is present in $\binom{m-1}{w-1}$ of the sets. Consequently, each component has $\binom{m-1}{w-1}$ occurrences. Performing this substitution results in
\begin{equation}
    \label{eq:proof_joint_like_marginal}
    \eta \binom{m-1}{w-1} \sum_{j=1}^{m} \pmf_{\RV{K}}(j \given \tparam{\sysparam}) = 1\,.
\end{equation}
Observe that $\sum_{j=1}^{m} \pmf_{\RV{K}}(j \given \tparam{\sysparam})$ must be equal to unity since we are summing over the entire sample space of $\RV{K}$. Performing this substitution and solving for $\eta$ yields
\begin{equation}
    \eta = \frac{1}{\binom{m-1}{w-1}}\,.
\end{equation}
\end{proof}

\begin{corollary}
\label{def:general_cond_cand_probability}
The conditional probability mass function $\pmf_{\rvcand \given \sv, \RV{W}}(\cand \given t, w, \tparam\sysparam)$ maps each candidate set $\cand$ to the probability that the given candidate set contains the component responsible for the system failure at time $t$. By the axioms of probability,
\begin{equation}
\label{eq:general_cond_cand_probability}
    \pmf_{\rvcand \given \sv, \RV{W}}(\cand \given t, w, \tparam\sysparam) = \frac
    {
        \sum_{j \in \cand} \haz_j\left(t \given \tparam\sysparam_j\right)
    }
    {
        \haz_{\sv}\left(t \given \tparam\sysparam\right)
    } \SetIndicator{\candsw{w}}(\cand)\,,
\end{equation}
where $\haz_{\sv}$ and $\haz_j$ are the failure rate functions of the system 
and the \jth component respectively.
\end{corollary}
\begin{proof}
By the axioms of probability,
\begin{equation}
    \pmf_{\rvcand \given \sv, \RV{W}}(\cand \given t, w, \tparam\sysparam) = \frac
    {
        \jointpdf(\cand, t \given w, \tparam\sysparam)
    }
    {
        \spdf(t \given \tparam\sysparam)
    }\,,
\end{equation}
where $\jointpdf(\; \cdot \given \tparam\sysparam)$ is the joint probability density function given by \cref{thm:general_joint_likelihood_cand}. Making this substitution immediately yields the result.
\end{proof}

\end{document}